\documentclass[utf8, russian]{eskdtext}
\usepackage{paralist}
% без определения что-то не работает
\newcommand{\No}{\textnumero}

\usepackage{titlesec} 
\titleformat{\section}
{\normalsize\bfseries} 
{\thesection} {1em}{} 
\titleformat{\subsection}
{\normalsize\bfseries}
{\thesubsection} {1em}{} 
\titleformat{\subsubsection}
{\normalsize\bfseries}
{\thesubsubsection} {1em}{}




% для штампа
\ESKDtitle{{\small Разработка MCP-сервера для Яндекс Трекера}}
\ESKDsignature{ТПЖА.090301.842 ПЗ}
\ESKDgroup{\footnotesize Кафедра ЭВМ\\ИВТб-4301-04-01}
\ESKDauthor{Бикметов}
\ESKDchecker{Коржавина}
%\ESKDnormContr{\resizebox{2.22cm}{\height}{Фамилия И.О.}}
%\ESKDapprovedBy{Фамилия И.О.}

% оступы от заголовков разделов и подразделов
%\ESKDsectSkip{section}{7mm}{7mm}
%\ESKDsectSkip{subsection}{5mm}{5mm}
%\ESKDsectSkip{subsubsection}{3mm}{3mm}

%%% Каждый раздел (\section) с новой страницы
\let\stdsection\section
\renewcommand \section{\newpage\stdsection}

\title{{\small Разработка MCP-сервера для Яндекс Трекера}}

\usepackage{enumitem}
%\setlist[enumerate,itemize]{leftmargin=0pt,itemindent=2cm}
\renewcommand{\labelitemi}{\normalfont\bfseries{--}}
\setlist{nosep,
left=0pt,
labelindent=\parindent,leftmargin=*,
leftmargin=0pt,itemindent=\labelwidth
}
\setlist[enumerate,2]{
leftmargin=\parindent,itemindent=\labelwidth
}
\setenumerate{listparindent=\parindent}
\makeatletter
\AddEnumerateCounter{\asbuk}{\@asbuk}{м)}
\makeatother
\renewcommand{\labelenumi}{\asbuk{enumi})} 
\renewcommand{\labelenumii}{\arabic{enumii})}

%%% Выравнивание и переносы %%%

\sloppy                             % Избавляемся от переполнений

\clubpenalty=10000                  % Запрещаем разрыв страницы после первой строки абзаца

\widowpenalty=10000                 % Запрещаем разрыв страницы после последней строки абзаца
% подключение частей документа

%Для кода
\usepackage{listings}
\usepackage{verbatim}

%Для ссылок
\usepackage{hyperref}

\begin{document}
\maketitle
%%% Делаем содержание
\tableofcontents
%%% Введение пишется без цифры и добавляется в содержание
\section*{Введение}
\addcontentsline{toc}{section}{Введение}

В современном мире разработка программного обеспечения становится все более сложной и требует интеграции различных инструментов и сервисов. Одной из ключевых задач является обеспечение эффективного взаимодействия между различными системами управления проектами и интеллектуальными агентами на основе Large Language Models (LLM). Яндекс Трекер представляет собой мощный инструмент для управления задачами и проектами, но для его полноценной интеграции с LLM-агентами необходимы специализированные решения.

MCP (Model Context Protocol) — это протокол, разработанный для предоставления информации о проекте различным инструментам и ассистентам. MCP-серверы обеспечивают доступ к модели проекта, позволяя интеллектуальным агентам понимать структуру и содержание проекта, что значительно упрощает автоматизацию и интеграцию.

Целью данной работы является разработка MCP-сервера для Яндекс Трекера, который будет обеспечивать интеграцию между системой управления задачами и LLM-агентами. Это позволит интеллектуальным агентам получать актуальную информацию о задачах, проектах и других сущностях Яндекс Трекера, улучшая их способность к автоматизации и принятию решений.

\newpage

\section{Техническое задание}
Техническое задание является основным документом, который описывает цели, задачи и требования к разработке проекта. В данной версии ТЗ  сформулированы ключевые аспекты, необходимые для успешного выполнения проекта, учитывая потребности всех заинтересованных сторон.

\subsection{Пользовательские истории}

\begin{enumerate}
    \item \textbf{Руководитель проекта:}
    \begin{itemize}
        \item Как руководитель проекта, я хочу иметь возможность отслеживать прогресс выполнения задач в Яндекс Трекере через LLM-ассистента, чтобы оперативно принимать управленческие решения и управлять ресурсами команды;
        \item Как руководитель проекта, я хочу, чтобы LLM-ассистент мог анализировать текущую загрузку команды и предлагать оптимизацию распределения задач, чтобы повысить эффективность работы и избежать перегрузок;
        \item Как руководитель проекта, я хочу получать автоматические отчеты о состоянии проекта и ключевых метриках через интеграцию с LLM-агентом, чтобы иметь актуальную информацию для презентаций и планирования.
    \end{itemize}
    
    \item \textbf{Продукт-менеджер:}
    \begin{itemize}
        \item Как продукт-менеджер, я хочу, чтобы LLM-ассистент мог анализировать бэклог задач и предлагать приоритизацию на основе бизнес-ценности и пользовательских требований, чтобы оптимизировать развитие продукта;
        \item Как продукт-менеджер, я хочу иметь возможность быстро создавать и обновлять пользовательские истории и требования через взаимодействие с LLM-агентом;
        \item Как продукт-менеджер, я хочу, чтобы LLM-ассистент мог автоматически формировать отчеты о выполнении эпиков и фич, чтобы отслеживать прогресс реализации стратегических целей продукта.
    \end{itemize}
    
    \item \textbf{Разработчик:}
    \begin{itemize}
        \item Как разработчик, я хочу, чтобы LLM-ассистент мог предоставлять контекст задачи из Яндекс Трекера непосредственно в моей IDE, чтобы быстрее понимать требования и начинать работу над задачей;
        \item Как разработчик, я хочу, чтобы LLM-агент мог генерировать шаблоны кода и документации на основе спецификаций задач в Яндекс Трекере, чтобы ускорить процесс разработки;
        \item Как разработчик, я хочу иметь возможность обновлять статус задачи и добавлять комментарии к задаче через LLM-ассистент, чтобы минимизировать переключение между инструментами и сосредоточиться на коде.
    \end{itemize}
    
    \item \textbf{Тестировщик:}
    \begin{itemize}
        \item Как тестировщик, я хочу, чтобы LLM-ассистент мог автоматически создавать тест-кейсы на основе описания задачи в Яндекс Трекере, чтобы ускорить процесс тестирования и повысить его качество;
        \item Как тестировщик, я хочу иметь возможность быстро добавлять баг-репорты в Яндекс Трекер через LLM-агента, интегрированный с моей IDE, чтобы сразу фиксировать найденные ошибки без переключения контекста;
        \item Как тестировщик, я хочу, чтобы LLM-ассистент мог анализировать историю изменений задачи и предлагать сценарии регрессионного тестирования, чтобы обеспечить полноту покрытия тестами после изменений в коде.
    \end{itemize}
\end{enumerate}


\subsection{Пользовательские сценарии}

Пользовательские сценарии описывают типовые ситуации взаимодействия различных ролей с MCP-сервером для Яндекс Трекера, интегрированным с LLM-агентами. Они охватывают ключевые аспекты работы с системой управления задачами через интеллектуальных ассистентов.

\begin{enumerate}
    \item \textbf{Сценарий 1. Отслеживание прогресса проекта руководителем.}\\
    \begin{itemize}
        \item Действующие лица: Руководитель проекта, LLM-ассистент, MCP-сервер, Яндекс Трекер;
        \item Цель: Получить актуальную информацию о состоянии проекта и ключевых метриках.
    \end{itemize}
    \begin{itemize}
        \item Руководитель проекта обращается к LLM-ассистенту с запросом "Покажи прогресс по проекту";
        \item LLM-ассистент отправляет запрос к MCP-серверу для получения данных из Яндекс Трекера;
        \item MCP-сервер извлекает информацию о задачах, статусах, сроках и загрузке команды;
        \item LLM-ассистент анализирует полученные данные и формирует отчет в понятной форме;
        \item Руководитель получает структурированную информацию о прогрессе проекта.
    \end{itemize}

    \item \textbf{Сценарий 2. Приоритизация задач продукт-менеджером.}\\
    \begin{itemize}
        \item Действующие лица: Продукт-менеджер, LLM-ассистент, MCP-сервер, Яндекс Трекер;
        \item Цель: Определить приоритеты задач на основе бизнес-ценности и пользовательских требований.
    \end{itemize}
    \begin{itemize}
        \item Продукт-менеджер запрашивает у LLM-ассистента анализ бэклога задач;
        \item LLM-ассистент обращается к MCP-серверу за данными о задачах, их описании, метках и комментариях;
        \item MCP-сервер предоставляет структурированные данные из Яндекс Трекера;
        \item LLM-ассистент анализирует информацию и предлагает приоритезацию на основе бизнес-ценности;
        \item Продукт-менеджер получает рекомендации и может внести корректировки через LLM-ассистент.
    \end{itemize}

    \item \textbf{Сценарий 3. Получение контекста задачи разработчиком.}\\
    \begin{itemize}
        \item Действующие лица: Разработчик, LLM-ассистент, MCP-сервер, Яндекс Трекер;
        \item Цель: Быстро получить полный контекст задачи для начала работы.
    \end{itemize}
    \begin{itemize}
        \item Разработчик начинает работу над задачей в своей IDE и запрашивает контекст;
        \item LLM-ассистент, интегрированный с IDE, обращается к MCP-серверу;
        \item MCP-сервер извлекает из Яндекс Трекера описание задачи, требования, комментарии и историю изменений;
        \item LLM-ассистент обрабатывает информацию и предоставляет разработчику структурированный контекст;
        \item Разработчик получает всю необходимую информацию для начала работы без переключения между инструментами.
    \end{itemize}

    \item \textbf{Сценарий 4. Создание баг-репорта тестировщиком.}\\
    \begin{itemize}
        \item Действующие лица: Тестировщик, LLM-ассистент, MCP-сервер, Яндекс Трекер;
        \item Цель: Быстро зафиксировать найденную ошибку в системе управления задачами.
    \end{itemize}
    \begin{itemize}
        \item Тестировщик находит баг и сообщает об этом LLM-ассистенту;
        \item LLM-ассистент собирает информацию о баге, включая шаги воспроизведения и ожидаемое поведение;
        \item LLM-ассистент отправляет запрос к MCP-серверу для создания задачи в Яндекс Трекере;
        \item MCP-сервер создает баг-репорт с заполненными полями на основе информации от LLM-ассистента;
        \item Тестировщик получает подтверждение создания баг-репорта и при необходимости может внести дополнения.
    \end{itemize}
\end{enumerate}


\subsection{Основания для разработки}
Основанием для разработки является задание на курсовой проект. Задание утверждено ООО «Философт», именуемый в дальнейшем Заказчиком, и Бикметовым Игорем Руслановичем (самозанятый), именуемым в дальнейшем исполнителем, Исполнитель обязан разработать «MCP-сервер для Яндекс Трекера» не позднее 26.12.2025. Наименование темы разработки – «Разработка MCP-сервера для Яндекс Трекера».

\subsection{Назначение разработки}
Данная система будет использоваться в рамках внутренней работы ООО «Философт» для интеграции Яндекс Трекера с LLM-агентами.

\subsubsection{Функциональные требования}

Ниже перечислены функциональные требования к MCP-серверу для Яндекс Трекера, сгруппированные по категориям:

\text{1. Управление задачами:}
\begin{itemize}
    \item Создание новых задач в системе Яндекс Трекер;
    \item Чтение информации о задачах по их идентификаторам;
    \item Обновление существующих задач с возвратом обновленной информации;
    \item Удаление задач;
    \item Поиск задач по различным критериям.
\end{itemize}

\text{2. Комментарии:}
\begin{itemize}
    \item Создание комментариев к задачам;
    \item Чтение комментариев задачи;
    \item Обновление комментариев;
    \item Удаление комментариев.
\end{itemize}

\text{3. Спринты:}
\begin{itemize}
    \item Создание новых спринтов;
    \item Чтение информации о спринтах;
    \item Обновление информации о спринтах;
    \item Удаление спринтов.
\end{itemize}

\text{4. Доски:}
\begin{itemize}
    \item Создание новых досок;
    \item Чтение информации о досках;
    \item Обновление информации о досках;
    \item Удаление досок.
\end{itemize}

\text{5. Статусы:}
\begin{itemize}
    \item Создание новых статусов;
    \item Чтение информации о статусах;
    \item Обновление информации о статусах;
    \item Удаление статусов.
\end{itemize}

\text{6. Очереди:}
\begin{itemize}
    \item Получение списка очередей задач;
    \item Получение детальной информации об очереди.
\end{itemize}

\text{7. Пользователи:}
\begin{itemize}
    \item Получение информации о текущем пользователе;
    \item Поиск пользователей по заданным критериям.
\end{itemize}

\text{8. Ресурсы (resources):}
\begin{itemize}
    \item Предоставление информации о доступных типах задач в организации;
    \item Предоставление списка приоритетов задач;
    \item Предоставление списка доступных статусов задач.
\end{itemize}

\text{9. Промпты (prompts):}
\begin{itemize}
    \item Анализ спринта или группы задач с предоставлением статистики и рекомендаций;
    \item Краткое изложение задачи с указанием её сути, статуса и исполнителей.
\end{itemize}

\subsubsection{Эксплуатационное назначение}
Система будет использоваться в рамках внутренней работы ООО «Философт» для оптимизации рабочего процесса.

\subsection{Требования к MCP-серверу}
В данном разделе будут описаны основные требования к MCP-серверу.
\subsubsection{Требования к составу выполняемых функций}

\begin{enumerate}
    \item \textbf{Создание задачи (create\_issue):}
    \begin{itemize}
        \item Входные данные: заголовок (summary), описание (description), очередь (queue), тип (type), приоритет (priority), исполнитель (assignee), компоненты (components), метки (labels);
        \item Выходные данные: объект созданной задачи с уникальным ключом, содержащий все основные атрибуты задачи.
    \end{itemize}
    
    \item \textbf{Чтение задачи (get\_issue):}
    \begin{itemize}
        \item Входные данные: идентификатор задачи (issueKey);
        \item Выходные данные: полная информация о задаче, включая статус, исполнителя, описание, комментарии, историю изменений.
    \end{itemize}
    
    \item \textbf{Обновление задачи (update\_issue):}
    \begin{itemize}
        \item Входные данные: идентификатор задачи (issueKey), поля для обновления (summary, description, status, assignee и др.);
        \item Выходные данные: обновленный объект задачи с актуальной информацией.
    \end{itemize}
    
    \item \textbf{Удаление задачи (delete\_issue):}
    \begin{itemize}
        \item Входные данные: идентификатор задачи (issueKey);
        \item Выходные данные: подтверждение успешного удаления задачи.
    \end{itemize}
    
    \item \textbf{Поиск задач (search\_issues):}
    \begin{itemize}
        \item Входные данные: критерии поиска (язык запросов Трекера), параметры сортировки, лимит и смещение результатов;
        \item Выходные данные: список найденных задач с основной информацией.
    \end{itemize}
    
    \item \textbf{Создание комментария (add\_comment):}
    \begin{itemize}
        \item Входные данные: идентификатор задачи (issueKey), текст комментария, упомянутые пользователи;
        \item Выходные данные: объект созданного комментария с идентификатором и содержимым.
    \end{itemize}
    
    \item \textbf{Чтение комментариев (get\_comments):}
    \begin{itemize}
        \item Входные данные: идентификатор задачи (issueKey), параметры пагинации;
        \item Выходные данные: список комментариев к задаче с информацией об авторе, дате и содержимом.
    \end{itemize}
    
    \item \textbf{Обновление комментария (update\_comment):}
    \begin{itemize}
        \item Входные данные: идентификатор комментария, новое содержимое;
        \item Выходные данные: обновленный объект комментария.
    \end{itemize}
    
    \item \textbf{Удаление комментария (delete\_comment):}
    \begin{itemize}
        \item Входные данные: идентификатор комментария;
        \item Выходные данные: подтверждение успешного удаления комментария.
    \end{itemize}
    
    \item \textbf{Создание спринта (create\_sprint):}
    \begin{itemize}
        \item Входные данные: название спринта, дата начала, дата окончания, доска, к которой привязан спринт;
        \item Выходные данные: объект созданного спринта с идентификатором и основными атрибутами.
    \end{itemize}
    
    \item \textbf{Чтение информации о спринте (get\_sprint):}
    \begin{itemize}
        \item Входные данные: идентификатор спринта;
        \item Выходные данные: полная информация о спринте, включая статус, даты, задачи, участников.
    \end{itemize}
    
    \item \textbf{Обновление спринта (update\_sprint):}
    \begin{itemize}
        \item Входные данные: идентификатор спринта, поля для обновления (даты, статус, название);
        \item Выходные данные: обновленный объект спринта.
    \end{itemize}
    
    \item \textbf{Удаление спринта (delete\_sprint):}
    \begin{itemize}
        \item Входные данные: идентификатор спринта;
        \item Выходные данные: подтверждение успешного удаления спринта.
    \end{itemize}
    
    \item \textbf{Создание доски (create\_board):}
    \begin{itemize}
        \item Входные данные: название доски, тип доски, фильтр задач, настройки отображения;
        \item Выходные данные: объект созданной доски с идентификатором и конфигурацией.
    \end{itemize}
    
    \item \textbf{Чтение информации о доске (get\_board):}
    \begin{itemize}
        \item Входные данные: идентификатор доски;
        \item Выходные данные: информация о доске, включая колонки, фильтры, настройки отображения.
    \end{itemize}
    
    \item \textbf{Обновление доски (update\_board):}
    \begin{itemize}
        \item Входные данные: идентификатор доски, поля для обновления (название, фильтры, настройки);
        \item Выходные данные: обновленный объект доски.
    \end{itemize}
    
    \item \textbf{Удаление доски (delete\_board):}
    \begin{itemize}
        \item Входные данные: идентификатор доски;
        \item Выходные данные: подтверждение успешного удаления доски.
    \end{itemize}
    
    \item \textbf{Создание статуса (create\_status):}
    \begin{itemize}
        \item Входные данные: название статуса, тип статуса (открыт, в работе, закрыт), цветовое обозначение;
        \item Выходные данные: объект созданного статуса с идентификатором и атрибутами.
    \end{itemize}
    
    \item \textbf{Чтение информации о статусе (get\_status):}
    \begin{itemize}
        \item Входные данные: идентификатор статуса;
        \item Выходные данные: информация о статусе, включая его тип, цвет, описание.
    \end{itemize}
    
    \item \textbf{Обновление статуса (update\_status):}
    \begin{itemize}
        \item Входные данные: идентификатор статуса, поля для обновления (название, тип, цвет);
        \item Выходные данные: обновленный объект статуса.
    \end{itemize}
    
    \item \textbf{Удаление статуса (delete\_status):}
    \begin{itemize}
        \item Входные данные: идентификатор статуса;
        \item Выходные данные: подтверждение успешного удаления статуса.
    \end{itemize}
    
    \item \textbf{Получение очередей (get\_queues):}
    \begin{itemize}
        \item Входные данные: параметры фильтрации и расширения;
        \item Выходные данные: список очередей с информацией о каждой очереди.
    \end{itemize}
    
    \item \textbf{Получение информации о очереди (get\_queue):}
    \begin{itemize}
        \item Входные данные: идентификатор очереди, параметры расширения;
        \item Выходные данные: детальная информация об очереди, включая настройки и права доступа.
    \end{itemize}
    
    \item \textbf{Получение информации о пользователе (get\_myself):}
    \begin{itemize}
        \item Входные данные: отсутствуют;
        \item Выходные данные: информация о текущем пользователе, включая права доступа.
    \end{itemize}
    
    \item \textbf{Поиск пользователей (search\_users):}
    \begin{itemize}
        \item Входные данные: строка поиска, лимит результатов;
        \item Выходные данные: список найденных пользователей с основной информацией.
    \end{itemize}
\end{enumerate}

\subsubsection{Требования к организации входных и выходных данных}

Входные и выходные данные MCP-сервера должны соответствовать спецификации Model Context Protocol (MCP) версии  05.11.2024. Все данные передаются в формате JSON.

\textbf{Входные данные:}
\begin{itemize}
    \item Все входящие запросы должны содержать корректный заголовок аутентификации (OAuth токен);
    \item Параметры запросов должны проходить валидацию с использованием схем, определенных в формате JSON Schema;
    \item Данные о задачах, пользователях, очередях и других сущностях должны соответствовать структуре, определенной в API Яндекс Трекера;
    \item Все строковые значения должны быть в кодировке UTF-8;
    \item Числовые значения должны соответствовать допустимым диапазонам, определенным в спецификации;
    \item Дата и время должны передаваться в формате ISO 8601.
\end{itemize}

\textbf{Выходные данные:}
\begin{itemize}
    \item Все ответы сервера должны содержать корректный HTTP-статус в зависимости от результата операции;
    \item В случае ошибки, сервер должен возвращать детализированное сообщение об ошибке в стандартном формате MCP;
    \item Все объекты задач, пользователей, очередей и других сущностей должны содержать полный набор атрибутов в соответствии с API Яндекс Трекера;
    \item Объем возвращаемых данных должен быть ограничен параметрами пагинации, если это применимо;
    \item Все данные должны быть защищены от несанкционированного доступа, токены и конфиденциальная информация не должны включаться в ответы сервера.
\end{itemize}

\subsubsection{Требования к временным характеристикам}
Максимальное время отклика системы 1 секунда.

\subsubsection{Требования к надежности}
Вероятность безотказной работы сервера должна составлять не менее 90.99\% при условии
исправности соединения с Яндекс Трекером.

\subsubsection{Требования к обеспечению надежного функционирования
системы}
Для обеспечения надежного функционирования системы необходимо выполнение следующих требований:
\begin{enumerate}
    \item Система должна быть спроектирована с учетом отказоустойчивости, чтобы минимизировать влияние возможных сбоев компонентов на общую работоспособность системы;
    \item Система должна поддерживать механизм мониторинга состояния компонентов и логирования событий для своевременного выявления и устранения неполадок;
    \item Должна быть обеспечена защита от несанкционированного доступа к данным и функциям системы с использованием современных методов аутентификации и авторизации;
    \item Система должна поддерживать масштабирование для обеспечения стабильной работы при увеличении нагрузки.
\end{enumerate}


\subsubsection{Время восстановления после отказа}
Время восстановления после отказа, вызванного неисправностью технических средств, не должно превышать времени,
требуемого на устранение неисправностей технических средств.

\subsubsection{Условия эксплуатации}

Для обеспечения надежной и безопасной эксплуатации MCP-сервера необходимо соблюдение следующих условий:

\begin{enumerate}
    \item MCP-сервер должен эксплуатироваться в закрытых помещениях с контролируемым климатом, без прямого воздействия солнечных лучей и атмосферных осадков;
    \item Температурный режим эксплуатации: от +10°C до +35°C;
    \item Напряжение питания: 220 В ±10\%, частота 50 Гц ±2\%;
    \item MCP-сервер должен быть защищен от пыли, вибраций и механических воздействий;
    \item Запрещена эксплуатация в условиях агрессивной среды (наличие химически активных газов, паров кислот и щелочей);
    \item Обслуживание MCP-сервера должно производиться квалифицированным персоналом, прошедшим инструктаж по технике безопасности;
    \item MCP-сервер должен быть подключен к стабилизированному источнику питания с возможностью резервного питания на время кратковременных отключений;
    \item Для MCP-сервера должна быть обеспечена надежная связь с Яндекс Трекером через защищенный канал связи;
    \item MCP-сервер должен находиться в контролируемом (дата-центре или серверной комнате) с ограниченным физическим доступом.
\end{enumerate}

\subsubsection{Климатические условия эксплуатации}
Специальные условия не требуются.

\subsubsection{Требования к видам обслуживания}

Обслуживание MCP-сервера включает в себя следующие виды работ:

\begin{enumerate}
    \item \text{Техническое обслуживание:}
    \begin{itemize}
        \item Регулярный контроль работоспособности системы (ежедневно);
        \item Проверка состояния логов и мониторинг производительности (ежедневно);
        \item Обновление программного обеспечения (по мере необходимости);
        \item Проверка целостности данных и резервных копий (еженедельно).
    \end{itemize}
    
    \item \text{Профилактическое обслуживание:}
    \begin{itemize}
        \item Проверка и тестирование резервного копирования данных (ежемесячно);
        \item Анализ производительности и выявление узких мест (ежемесячно);
        \item Проверка конфигурации безопасности и настройка доступа (ежеквартально);
        \item Проверка и обновление документации по системе (ежеквартально).
    \end{itemize}
    
    \item \text{Ремонтное обслуживание:}
    \begin{itemize}
        \item Устранение неисправностей в работе системы (по мере возникновения);
        \item Восстановление работоспособности после сбоев (при необходимости);
        \item Замена или восстановление компонентов системы (при необходимости);
        \item Восстановление из резервной копии в случае потери данных (при необходимости).
    \end{itemize}
\end{enumerate}

Все виды обслуживания должны выполняться квалифицированным персоналом с соблюдением установленных процедур и регламентов безопасности.

\subsubsection{Требования к численности и квалификации персонала}

Обслуживание MCP-сервера должно осуществляться специалистом с высшим техническим образованием и опытом работы в области разработки и сопровождения серверных приложений не менее 2 лет. Для обеспечения бесперебойной работы системы рекомендуется наличие дежурного системного администратора с соответствующей квалификацией.

\subsubsection{Требование к маркировке и упаковке}
Специальных требований к маркировке и упаковке не предъявляется.

\subsubsection{Требования к транспортированию и хранению}
Специальных требований не предъявляется.

\subsubsection{Специальные требования}
Специальных требований не предъявляется.

\subsection{Требования к программной документации}
Предварительный состав программной документации:
\begin{enumerate}
    \item техническое задание (включает описание применения);
    \item методика испытаний;
    \item документация по MCP-серверу для Яндекс Трекера.
\end{enumerate}
    
\subsection{Технико-экономические показатели}
MCP-сервер будет использоваться в рамках внутренней работы ООО «Философт». Это позволит оптимизировать рабочий процесс персонала. Что, в свою очередь, сократит время на реализацию поставленных задач.

\subsection{Стадии и этапы разработки}
Разработка должна быть проведена в три стадии:
\begin{enumerate}
    \item техническое задание;
    \item технический (и рабочий) проекты;
    \item внедрение.
\end{enumerate}
На стадии «Техническое задание» должен быть выполнен этап разработки, согласования и
утверждения настоящего технического задания.
На стадии «Технический (и рабочий) проект» должны быть выполнены перечисленные
ниже этапы работ:
\begin{enumerate}
    \item разработка сервера;
    \item разработка документации;
    \item испытания сервера.
\end{enumerate}
На стадии «Внедрение» должен быть выполнен этап разработки «Подготовка и передача
сервера».
Содержание работ по этапам:
На этапе разработки технического задания должны быть выполнены перечисленные ниже
работы:
\begin{enumerate}
    \item постановка задачи;
    \item определение и уточнение требований к техническим средствам;
    \item определение требований к серверу;
    \item определение стадий, этапов и сроков разработки сервера и документации на него;
    \item согласование и утверждение технического задания.
\end{enumerate}

\subsection{Порядок контроля и приемки}
Приемосдаточные испытания MCP-сервера должны проводиться согласно разработанной
исполнителем и согласованной заказчиком «Методики испытаний».
Ход проведения приемо-сдаточных испытаний заказчик и исполнитель документируют в
протоколе испытаний.
На основании протокола испытаний исполнитель совместно с заказчиком подписывают акт
приемки-сдачи MCP-сервера в эксплуатацию.
\subsection{Заключение}

Настоящее техническое задание определяет требования к разработке MCP-сервера для Яндекс Трекера, предназначенного для интеграции с LLM-агентами. Документ охватывает все аспекты разработки, включая функциональные и нефункциональные требования, условия эксплуатации, требования к персоналу и документации. Выполнение всех указанных в техническом задании требований обеспечит создание надежного и эффективного решения для автоматизации работы с Яндекс Трекером. Техническое задание является основанием для выполнения работ по разработке и последующему внедрению MCP-сервера в ООО «Философт».

\section{Анализ предметной области}

\subsection{Общая характеристика системы}

MCP-сервер для Яндекс Трекера представляет собой специализированное программное решение, предназначенное для обеспечения интеграции между системой управления задачами Яндекс Трекер и интеллектуальными агентами на основе Large Language Models (LLM). Основной целью системы является предоставление LLM-агентам доступа к информации о проекте через протокол Model Context Protocol (MCP), что позволяет автоматизировать взаимодействие с задачами, улучшить контекстное понимание проектных сущностей и оптимизировать рабочие процессы в командной разработке.

Система функционирует как посредник между Яндекс Трекером и LLM-ассистентами, обеспечивая выполнение широкого спектра операций: создание, чтение, обновление и удаление задач, комментариев, спринтов, досок и других элементов проекта. Также сервер предоставляет информацию о пользователях, очередях задач, статусах и других ресурсах, необходимых для полноценной интеграции. Важной особенностью является поддержка промптинга - возможность анализа данных проекта и генерации рекомендаций на основе искусственного интеллекта.

Сервер должен соответствовать требованиям производительности (максимальное время отклика 1 секунда) и надежности (вероятность безотказной работы не менее 90.99\%), а также обеспечивать безопасное взаимодействие с внешней системой через аутентификацию с использованием OAuth токенов. Архитектура сервера должна быть спроектирована с учетом отказоустойчивости и возможности масштабирования при увеличении нагрузки.

\subsection{Анализ существующих проблем}

В современных условиях разработка программного обеспечения требует слаженного взаимодействия между различными ролями в команде: руководителями, продукт-менеджерами, разработчиками и тестировщиками. В организации применяется Яндекс Трекер для управления проектами и часто возникают следующие проблемы:

\begin{enumerate}
    \item \textbf{Фрагментация информации}: Команда разработки сталкивается с необходимостью переключения между различными инструментами (IDE, система управления задачами, чаты и др.) для выполнения повседневных задач. Это приводит к снижению продуктивности и увеличению времени на выполнение рутинных операций;
    
    \item \textbf{Недостаточная автоматизация}: Рутинные операции, такие как создание задач, обновление статусов, добавление комментариев, требуют ручного вмешательства и отнимают значительное количество времени у специалистов;
    
    \item \textbf{Сложности в анализе данных}: Руководителям и продукт-менеджерам сложно оперативно получать аналитическую информацию о состоянии проекта, загрузке команды и приоритетности задач без ручного анализа данных в Яндекс Трекере;
    
    \item \textbf{Отсутствие интеллектуальной поддержки}: Отсутствует интеграция с интеллектуальными агентами, которые могли бы автоматизировать анализ задач, генерировать отчеты, предлагать оптимизацию распределения задач и создавать тест-кейсы;
    
    \item \textbf{Задержки в коммуникации}: Коммуникация между членами команды может замедляться из-за необходимости переключения между инструментами для получения контекста задачи, особенно для разработчиков, которым нужно понимать требования и спецификации задачи;
    
    \item \textbf{Недостаточная интеграция с ИИ}: Отсутствует прямая интеграция Яндекс Трекера с LLM-агентами, что не позволяет использовать возможности искусственного интеллекта для автоматизации процессов, анализа данных и принятия решений.
\end{enumerate}

Данные проблемы особенно актуальны для команд, которые используют Яндекс Трекер как основной инструмент управления проектами, но не имеют полноценной интеграции с современными ИИ-ассистентами. Решение этих проблем возможно через разработку MCP-сервера, который обеспечит прямое взаимодействие между Яндекс Трекером и LLM-агентами, позволяя автоматизировать рутинные задачи и улучшить общую эффективность работы команды.

\begin{center}
\text{Таблица 1 - Анализ существующих проблем}
\end{center}

\begin{center}
\begin{tabular}{|p{8cm}|p{8cm}|}
\hline
\textbf{Название проблемы} & \textbf{Краткое описание} \\
\hline
Фрагментация информации & Переключение между инструментами снижает продуктивность \\
\hline
Недостаточная автоматизация & Рутинные операции требуют ручного вмешательства \\
\hline
Сложности в анализе данных & Трудности с получением аналитической информации \\
\hline
Отсутствие интеллектуальной поддержки & Нет интеграции с ИИ-агентами \\
\hline
Задержки в коммуникации & Затруднено получение контекста задачи \\
\hline
Недостаточная интеграция с ИИ & Отсутствует прямая интеграция с LLM-агентами \\
\hline
\end{tabular}
\end{center}

\subsection{Анализ требований по решению существующих проблем}

С точки зрения системного аналитиза и заказчика, MCP-сервер должен решать выявленные проблемы за счет реализации следующих требований и функциональных возможностей:

\textbf{1. Устранение фрагментации информации:}
\begin{itemize}
    \item Обеспечение прямой интеграции с IDE и инструментами разработки через MCP-протокол, что позволяет получать информацию из Яндекс Трекера без переключения между приложениями;
    \item Предоставление контекста задачи непосредственно в рабочем окружении разработчика;
    \item Генерация краткого изложения задачи с указанием её сути, статуса исполнителей;
    \item Обогащение контекста LLM-модели под реализацию конкретной задачи.
\end{itemize}

\textbf{2. Повышение уровня автоматизации:}
\begin{itemize}
    \item Реализация полного цикла операций с задачами, комментариями, спринтами и досками;
    \item Автоматизация создания и обновления задач через интеграцию с LLM-ассистентами, что исключает необходимость ручного ввода данных в интерфейсе Яндекс Трекера;
    \item Поддержка поиска задач с различными критериями для автоматического анализа и обработки задач.
\end{itemize}

\textbf{3. Упрощение анализа данных:}
\begin{itemize}
    \item Делегирование анализа данных LLM-модели;
    \item Предоставление информации о пользователях и очередях для анализа загрузки команды;
    \item Анализ спринта или группы задач с предоставлением статистики и рекомендаций;
    \item Получение детальной информации о задачах, включая историю изменений и комментарии, для автоматического формирования отчетов.
\end{itemize}

\textbf{4. Обеспечение интеллектуальной поддержки:}
\begin{itemize}
    \item Интеграция с LLM-агентами через MCP-протокол для автоматического анализа задач и генерации рекомендаций;
    \item Генерация тест-кейсов на основе описания задачи через интеллектуальный анализ содержимого задач;
    \item Приоритезация задач на основе анализа бизнес-ценности и пользовательских требований.
\end{itemize}

\textbf{5. Ускорение коммуникации:}
\begin{itemize}
    \item Обеспечение мгновенного доступа к контексту задачи из IDE или других инструментов разработки без необходимости переключения на интерфейс Яндекс Трекера;
    \item Быстрое добавление комментариев и обновления статусов задач непосредственно из рабочего окружения разработчика;
    \item Интеграция системами чатов и коммуникации для автоматического уведомления участников проекта о изменениях.
\end{itemize}

\textbf{6. Интеграция с ИИ-технологиями:}
\begin{itemize}
    \item Обеспечение соответствия спецификации Model Context Protocol для полноценной интеграции с LLM-агентами;
    \item Предоставление информации о типах задач, приоритетах и статусах в формате, понятном для ИИ-ассистентов;
    \item Реализация безопасного и надежного взаимодействия с Яндекс Трекером через аутентификацию и защиту от несанкционированного доступа.
\end{itemize}

Таким образом, MCP-сервер решает существующие проблемы за счет создания единой точки интеграции между Яндекс Трекером и интеллектуальными агентами, обеспечивая автоматизацию рутинных операций, улучшение контекста для разработчиков и аналитиков, а также повышение общей эффективности командной работы. Выполнение всех указанных требований позволит значительно сократить время на реализацию проектов и повысить качество взаимодействия между членами команды.

\subsection{Обзор существующих решений}

В рамках анализа предметной области и постановки задачи разработки MCP-сервера для Яндекс Трекера, было выявлено несколько существующих аналогичных решений, которые обеспечивают интеграцию Яндекс Трекера с интеллектуальными агентами через MCP-протокол. Ниже представлены основные аналоги и их характеристики:

\textbf{1. Официальный MCP-сервер от Yandex Cloud:}
\begin{itemize}
    \item Официальный шаблон MCP-сервера для Yandex Tracker, представленный в Yandex AI Studio;
    \item Позволяет AI-агенту полноценно работать с задачами и другими сущностями Yandex Tracker;
    \item Поддерживаются инструменты получения информации о задаче, проекте, портфеле и цели;
    \item Реализованы следующие инструменты: получение связей задачи, создание задачи, комментария к задаче, цели, изменение статуса задачи, параметров задачи, параметров цели, массовое изменение задач, статусов задач, проектов, портфелей и целей, массовый перевод задач в другую очередь;
    \item Требует аутентификационные данные: токен доступа и идентификатор организации.
\end{itemize}

\textbf{2. MCP-сервер от пользователя urfv (GitHub репозиторий urfv/yandex-tracker-mcp):}
\begin{itemize}
    \item MCP-сервер для интеграции Яндекс Трекера с Claude через Model Context Protocol (MCP);
    \item Базовая версия с минимумом методов, которая позволяет Claude просматривать и управлять задачами в Яндекс Трекере;
    \item Поддерживает получение проекта, получение задачи, получение комментариев к задаче, создание задачи, редактирование задачи, перемещение задачи, подсчет задач, поиск задач, связывание задач;
    \item Использует Python в качестве языка реализации;
    \item Имеет лицензию MIT.
\end{itemize}

\textbf{3. MCP-сервер от пользователя aikts (GitHub репозиторий aikts/yandex-tracker-mcp):}
\begin{itemize}
    \item Комплексный Model Context Protocol (MCP) сервер, который позволяет AI-ассистентам взаимодействовать с Yandex Tracker API;
    \item Обеспечивает безопасный, аутентифицированный доступ к задачам, очередям, комментариям, учету рабочего времени и функциональности поиска Yandex Tracker с дополнительным кэшированием Redis для повышения производительности;
    \item Поддерживает полный набор функций: управление очередями, управление пользователями, операции с задачами, управление полями, управление статусами и типами, расширенный язык запросов, кэширование производительности, настраиваемые элементы управления безопасности, несколько транспортных опций, аутентификацию OAuth 2.0, поддержку организации;
    \item Поддерживает подключение к различным MCP-клиентам, включая Claude Desktop, Claude Code, Cursor, Windsurf, Zed, GitHub Copilot (VS Code) и другие;
    \item Использует Python в качестве языка реализации;
    \item Имеет лицензию Apache-2.0.
\end{itemize}

\textbf{Отличительные особенности разрабатываемого решения:}
\begin{itemize}
    \item Специализированная интеграция с Яндекс Трекером, адаптированная под требования конкретного проекта;
    \item Поддержка полного спектра функций MCP-протокола для взаимодействия с LLM-агентами;
    \item Возможность кастомизации под различные LLM-модели и сценарии использования;
    \item Открытая архитектура для расширения функциональности;
    \item Учет специфики внутреннего использования в ООО «Философт».
\end{itemize}

Таким образом, анализ существующих аналогов показывает, что, несмотря наличие нескольких решений для интеграции Яндекс Трекера с MCP-протоколом, разрабатываемое решение имеет свои уникальные особенности и преимущества, обусловленные спецификой требований заказчика и внутреннего использования.

\end{document}
